% !PRINTED 526
\weberpage{526}
%<<CONT
(by the adjunction of variables) is enlarged. There thus arises an extended field, in which the laws of decomposability hold purely.

The starting point of the definitions was formed by the field $R$ of the rational numbers, and if we wish to give ourselves an account of the question on which properties of the field $R$ the possibility of this extension rests, we find that it is on the one hand the existence of the integers in $R$, and on the other hand the unique decomposability of these integers into prime factors, that alone were used in the whole deduction. If, then, we let any other field to which these two properties belong take the place of the field $R$, we shall be able to draw the same conclusions. If we take for $R$ another algebraic number field, we obtain, to be sure, nothing new; but we certainly do if, for example, we put in place of the field $R$ the field of the rational functions of one variable, to which, indeed, the two fundamental properties also belong. In this way we gain a starting point for the theory of the algebraic functions of one variable\footnote{Cf. \gesperrt{Dedekind-Weber}, Theorie der algebraischen Functionen einer Veränderlichen. \gesperrt{Crelle's} Journal, Bd. 92.}.

\begin{center}
\rule{4cm}{0.4pt}
\end{center}
